% Generated draft: requires mathematical and source review.
\documentclass[11pt]{article}
\usepackage[margin=25mm]{geometry}
\usepackage{fontspec}
\setmainfont{Latin Modern Roman}
\usepackage{amsmath,amssymb,mathtools,unicode-math}
\setmathfont{Latin Modern Math}
\usepackage{enumitem,needspace,xurl,hyperref}
\setlength{\parindent}{0pt}
\setlength{\parskip}{5pt}
\setlength{\emergencystretch}{3em}
\newcommand{\N}{\mathbb{N}}
\newcommand{\Z}{\mathbb{Z}}
\newcommand{\Q}{\mathbb{Q}}
\newcommand{\R}{\mathbb{R}}
\newcommand{\C}{\mathbb{C}}
\newcommand{\F}{\mathbb{F}}
\newcommand{\A}{\mathbb{A}}
\newcommand{\PP}{\mathbb{P}}
\newcommand{\E}{\mathbb{E}}
\newcommand{\eps}{\varepsilon}
\newcommand{\dd}{\,\mathrm{d}}
\begin{document}
\section*{UnsolvedMath ID: 701315}
\textbf{Source problem code:} catalogue entry
\subsection{Definitions and mathematical statement}
Let \(p\) be a prime and let \(\Omega\) be an infinite set. For a permutation \(g\) of \(\Omega\), define its support by
\[
\operatorname{supp}(g)=\{\omega\in\Omega:g(\omega)\ne\omega\}.
\]
Write \(\operatorname{Sym}_{\mathrm{fin}}(\Omega)\) for the group of permutations with finite support. A subgroup \(G\leq\operatorname{Sym}_{\mathrm{fin}}(\Omega)\) is a \(p\)-group if every element has order a power of \(p\). A nonempty subset \(B\subseteq\Omega\) is a block for \(G\) if, for every \(g\in G\), either \(g(B)=B\) or \(g(B)\cap B=\varnothing\).

For a group \(H\) and \(h\in H\), write
\[
h^H=\{x^{-1}hx:x\in H\}.
\]
The group \(H\) is an \(FC\)-group if \(h^H\) is finite for every \(h\in H\). A minimal non-\(FC\)-subgroup of \(G\) means a subgroup \(H\leq G\) that is not an \(FC\)-group and whose every proper subgroup is an \(FC\)-group. Minimality here is with respect to subgroup inclusion.

The source also requires the permutation action of \(G\) on \(\Omega\) to be \emph{totally imprimitive}. The supplied record does not define this term. It is retained as an unresolved source qualification; the formulation below depends on its intended definition.

\textbf{Question (A. O. Asar, Kourovka Notebook Problem 17.13).} For every prime \(p\), every infinite set \(\Omega\), and every totally imprimitive \(p\)-group \(G\leq\operatorname{Sym}_{\mathrm{fin}}(\Omega)\), suppose that the support of every nontrivial cycle in the disjoint cycle decomposition of every \(g\in G\) is a block for \(G\). Must there exist \(H\leq G\) such that
\[
\bigl(\exists h\in H:\ |h^H|=\infty\bigr)
\quad\text{and}\quad
\bigl(\forall K<H\ \forall k\in K:\ |k^K|<\infty\bigr)?
\]
Here \(K<H\) means that \(K\) is a proper subgroup of \(H\). If fixed points are included as one-point cycles, their underlying singleton sets are automatically blocks, so they impose no additional condition.
\subsection{Short English statement}
Asar asks whether every totally imprimitive group of finitary permutations of prime-power orders, satisfying the stated block condition on all cycle supports, contains a subgroup that is not an \(FC\)-group although every proper subgroup is an \(FC\)-group. The intended definition of ``totally imprimitive'' remains to be confirmed from the source.
\subsection{Sources}
\begin{itemize}
\item Supplied local UnsolvedMath registry record: permanent ID \(701315\), exact source problem code \texttt{catalogue entry}, title \emph{Kourovka Notebook Problem 17.13}, upstream revision \texttt{d371ee917ac9647fbb06cd6fc529deec6f3dda82}. The registry labels the problem \texttt{open}; this is a supplied status, not an independently verified mathematical or literature assessment.
\item Attributed source: \emph{The Kourovka Notebook (Unsolved Problems in Group Theory)}, 17th issue (2010), Problem 17.13, proposed by A. O. Asar. The supplied background also identifies the 21st issue edition (2026), arXiv:1401.0300v48.
\item The supplied collection description credits E. I. Khukhro and V. D. Mazurov as editors of the Kourovka Notebook collection covering Issues 1--21 (1965--2026). These bibliographic details are supplied attribution and have not been independently checked.
\end{itemize}
\end{document}
